\documentclass[a4paper, headsepline]{scrartcl}
\usepackage[T1]{fontenc} % speziell für deutsch
\usepackage[english]{babel}
\usepackage[utf8]{inputenc}
\usepackage{hyperref}
\usepackage[pdftex]{graphicx}
\usepackage{tabularx}
\usepackage{geometry}
\usepackage{scrpage2}
\usepackage{graphicx}
\usepackage{listings}
\usepackage{color}
\usepackage{lstautogobble}
\usepackage{float}
\usepackage{microtype} %verhindert zu breite Texte
\usepackage{setspace}
\usepackage{abstract}
\usepackage[justification=justified,singlelinecheck=false]{caption}
\usepackage[normalem]{ulem}
\useunder{\uline}{\ul}{}
\usepackage[toc,page]{appendix}


%Zeilenabstand
\onehalfspacing


%\restylefloat{figure}

 
%Seitenrandeinstellungen
%\geometry{top=12mm,left=25mm,right=25mm,bottom=13mm, includeheadfoot, headsep=12mm, foot=1cm}

%kein Einrücken nach Absaetzen
\setlength{\parindent}{0em} 


%Kopfzeileneinstellungen
\clearscrheadfoot 
%\pagestyle{scrheadings}  %linie unter Kopfzeile
%\automark[subsection]{section} 
%\renewcommand{\sectionmark}[1]{\markboth{#1}{}}% 
%\ihead{Federated Identities for Worldwide Scientific Collaboration\hfill\leftmark}
%\setheadsepline{0.2pt}
\setlength{\headheight}{1.1\baselineskip}

%Fusszeileneinstellung
\cfoot{\pagemark}
 
%kein Mensch will Boxen um Inhaltsverzeichnispunkte
\hypersetup{%
    pdfborder = {0 0 0}
}

%Zitierweise
\usepackage[backend=biber, style=ieee, natbib=true, hyperref=false,firstinits=true, bibencoding=utf8]{biblatex}
\renewcommand*{\labelnamepunct}{\addcolon\addspace}
\renewcommand*{\newunitpunct}{\addcomma\addspace}
\renewcommand*{\finentrypunct}{}
\DeclareFieldFormat{title}{#1\isdot}

%erzwingt Blocksatz auch wenn dadurch Lücken zwischen Wörtern entstehen
%nur wenn mans nicht anders will :)
%\sloppy

% Befehl für Anführungszeichen %Stefan: Wie folgt anwenden: \qte{Dieser Text ist in anführungszeichen}. Okay? Keine normalen Anführungszeichen reinschreiben
\newcommand{\qte}[1]{\glqq#1\grqq}
% \qte {Hase}

%Auch Paragraphen nummerieren und ins Inhaltsverzeichnis aufnehmen
\setcounter{secnumdepth}{4}  % = Nummerierung vertiefen *
\setcounter{tocdepth}{4}   %= Aufnahme in das Inhaltsverzeichnis *

%% ################# Listings ###################

\usepackage[section]{minted}


\graphicspath{{Pictures/}}

\bibliography{literatur}
%\bibliographystyle{plain}
	
\usepackage[babel,english]{csquotes}
%Diese Felder sollen nicht im Inhaltsverzeichnis angzeigt werden
\AtEveryBibitem {
	\clearlist{address}
	\clearfield{date}
	\clearfield{eprint}
	\clearfield{isbn}
	\clearfield{issn}
	\clearlist{location}
	\clearfield{month}
	\clearfield{series}
	\clearfield{edition}
	\clearfield{pagetotal}}
	
%% Einbinden der bib-Datei (=Quellen)
\bibliography{Literatur.bib}
\begin{document}
	\pagenumbering{Roman}
	
	
	\begin{titlepage}
		\begin{center}
			\includegraphics[scale=1.5]{Deckblatt/HTW_Logo_rgb.jpg} \qquad \qquad \\[4ex]
			\Large{\textbf{Web Applications}} \\[8ex]
			\LARGE{\textbf{Watering of Things}}\\[3ex]
			\large{\quad} \\ %Schriftgröße fürs Deckblatt definieren (wenn diese Zeile entfernt wird ist alles weitere in LARGE
			\begin{tabular}{l l l} \\
				Verfasser: & \quad Doreen Sacker, Tobias Wochinger \\[2ex]
				Studiengang: & \quad Internationale Medieninformatik (Master)\\
                 & \quad Angewandte Informatik (Master) \\[2ex]
				%Fachbereich: & \quad Wirtschaftswissnschaften II \\[2ex]
				Matrikelnummern: & \quad 552936, 552880 \\[2ex]
				Zuständiger Prof.: & \quad Prof. Dr. Gefei Zhang \\[2ex]
				%& \quad Yves Kemp \\[2ex]
				Datum: & \quad 01. Februar 2017 %\quad (Ver. 1.03)
			\end{tabular}
		\end{center}
	\end{titlepage}
	\newpage
	\tableofcontents
	\newpage
	\pagenumbering{arabic}
    %\vspace*{5cm}
	% \begin{abstract}
	\label{Abstract}

\end{abstract}
    %\newpage
    \section{Einleitung}
Nachdem JavaScript früher überwiegend bei der Entwicklung im Frontend zum Einsatz kam, profitieren JavaScript-Entwickler immer mehr vom Paradigma \textit{Learn once, write anywhere}. So schickt sich JavaScript an, mit dem Boom von NodeJS auch serverseitig verstärkt eine Rolle zu spielen. Mobile Applikationen konnten jedoch bisher nur eingeschränkt mittels JavaScript umgesetzt werden. Zu groß waren die Limitierungen durch den Fakt, dass mit JavaScript erstellte Applikationen nur innerhalb einer Web-View auf den Geräten ausgeführt werden konnten. Auf einem Hackathon entwickelte Facebook jedoch 2013 eine Möglichkeit React-Komponenten durch native UI-Komponenten auf den Geräten zu visualisieren. Dieser Prototyp wurde daraufhin innerhalb interner Projekte bei Facebook weiterentwickelt, bis eine erste Version von React Native schließlich im März 2015 als Open-Source-Projekt veröffentlicht wurde. Facebook hatte bei der Entwicklung ursprünglich lediglich die iOS-Plattform im Blick. Rasch wurde jedoch die Möglichkeit der Wiederverwendung für Android erkannt, sodass React Native seit Sptember 2015 auch für Android verfügbar ist.
Die Bereitstellung von React Native als Open-Source-Projekt sorgte von Beginn an für eine große Community, die React Native seitdem unter Facebooks Schirmherrschaft in großen Schritten vorantreibt. Über 9000 Commits von mehr als 1000 Entwicklern sowie ein zweiwöchiger-Release-Rhythmus zeugen von dem Bedarf und dem Potenzial der  Entwicklung nativer Smartphone-Applikationen mittels JavaScript.
    \section{Grundlagen}

\subsection{React}
\textit{React} ist eine 2013 von Facebook vorgestellte Bibliothek zur programmatischen Erstellung von UI-Komponenten mittels JavaScript \cite{facebook_inc._react_2016}. Im Gegensatz zu Frameworks wie \textit{AngularJS} oder \textit{Ember.js} handelt sich hierbei nicht um die Umsetzung des MV*-Paradigmas. Stattdessen fokussiert sich React auf die View-Komponente. Während Web-Frameworks die View meist mittels Templates beziehungsweise HTML beschreiben, setzt React diese mittels JavaScript um \cite{hunt_why_2013}. Die Motivation von Facebook zur Entwicklung von React war es, Single-Page-Anwendungen zu schreiben, deren Komponenten sich bei Datenänderung selbst updaten. Um dabei die kostenintensiven Änderungen am DOM zu minimieren, arbeitet React mit einem sogenannten \textit{Virtual DOM}. Änderungen werden zuerst auf den Virtual DOM angewandt, welcher diese performant verarbeiten kann. Anhand eines Vergleichs werden die minimal durchzuführenden Änderungen festgestellt und auf den eigentliche DOM angewandt. Dies führt dazu, dass React eine sehr performante Möglichkeit ist Webanwendungen mit sich dynamisch ändernden Views umzusetzen. Weiterhin kann React sowohl Client- als auch Server-seitig gerendert werden \cite{hunt_why_2013}.
 
 %TODO oneDirectional
\subsection{JSX}
Die Erstellung von View-Komponenten mittels JavaScript und React wirkt insbesondere für HTML-gewöhnte Entwickler wenig intuitiv. Beispielsweise wird mit folgendem Code ein div-Element mit dem Text \textit{Hello} erzeugt.
\begin{listing}[H]
    \begin{minted}{js}
   const element = React.createElement(
        "div",
        null,
        "Hello"
        );
    }
    \end{minted}
    \caption{Erstellung eines div-Elements mit dem Text \textit{Hello} ohne JSX}
    \label{lst:jsx_without_jsx}
\end{listing}

Meist wird daher die Syntax-Erweiterung \textit{JSX} genutzt. JSX ist ein an HTML angelehnter \textit{Syntactic Sugar}, welcher es erlaubt Listing \ref{lst:jsx_without_jsx} wie folgt zu schreiben:
\begin{listing}[H]
    \begin{minted}{js}
const element = <div>Hello</div>;
    \end{minted}
    \caption{Erstellung eines div-Elements mit dem Text \textit{Hello} mit JSX}
    \label{lst:jsx_with_jsx}
\end{listing}
Diese neue Syntax muss vor Ausführung jedoch wieder zu konventionellem JavaScript übersetzt werden. Dies kann beispielsweise mit Babel geschehen \cite{facebook_inc._introducing_2016}.
Innerhalb der mittels JSX-Komponenten eingesetztes JavaScript muss innerhalb von geschweiften Klammern definiert werden: 

\begin{listing}[H]
    \begin{minted}{js}
const userName = 'Max';
const element =  <div>Hello {userName}</div>;
    \end{minted}
    \caption{Verwendung von JavaScript-Ausdrücken innerhalb von JSX-Komponenten}
    \label{lst:embedded_expression}
\end{listing}

Mehrere Elemente können wie in HTML geschachtelt werden, sofern sie sich in einem gemeinsamen Elternelement befinden.

\begin{listing}[H]
    \begin{minted}{js}
const userName = 'Max';
const element = (
    <div>
        <div>Hello</div>
        <div>{userName}</div>
    </div>
    );
    \end{minted}
    \caption{Geschachteltes JSX}
    \label{lst:jsx_nested}
\end{listing}

Oft ist es weiterhin notwendig Daten von den Elternkomponenten in die Kindkomponenten zu übergeben. Dies wird in JSX mittels sogenannter \textit{Properties} umgesetzt.

\begin{listing}[H]
    \begin{minted}{js}
class Hello extends React.Component {
    render() {
        return (
            <div>Hello {this.props.userName}</div>
        );
    }
}

ReactDOM.render(<Hello userName='Max' />, mountNodeInDOM);
    \end{minted}
    \caption{Übergabe von Properties an Kind-Komponenten}
    \label{lst:jsx_nested}
\end{listing}



% 
% Properties
% Unterlemente


\subsection{Cross-Platform-Mobile-Development}
Applikationen für für mobile Plattformen werden traditionell in drei Kategorien unterteilt. \textit{Native}, \textit{hybride} und \textit{Web-Apps}. Native Apps werden speziell für die jeweilige Plattform, z.B. iOS oder Android, und nutzen die APIs der Plattform direkt. Dadurch sind native Apps am Besten an die jeweilige Plattform angepasst und erzielen hohe Performance. Naturgemäß besitzen native Apps das Look-and-Feel der jeweiligen Plattform. Dieses Vorgehen hat jedoch den Nachteil, dass für jede Plattform eine eigenständige App entwickelt werden muss sowie spezifisches Know-How pro Plattform nötig ist. Daher verbreiten sich immer öfters HTML5-basierte Web-Apps. Dieses müssen nicht installiert werden, sondern werden im Browser des Gerätes ausgeführt. Somit können sie wie normale Websiten mit Technologien wie HTML5, CSS und JavaScript erstellt werden. Nachteil der Webapps ist die geringe Performance im Vergleich zu nativen Apps und dass sie nur auf einen Teil der Hardwareressourcen, wie z.B. GPS, zugreifen können. Den Kompromiss zwischen nativen und Web-Apps bilden hybride Apps. Hier wird letztlich eine Webview in eine native App gewrapped. Damit verhalten sich hybride Apps für den Nutzer wie native Apps. Hybride Apps können auf alle Hardwarefunktionen des Geräts nutzen. Da die View weiter im Browser gerendert wird, ist jedoch auch hier die Performance geringer. Bekannte Frameworks für hybride Apps sind beispielsweise PhoneGap \cite{adobe_systems_inc._adobe_2016} oder Cordova \cite{the_apache_software_foundation_apache_2016}.
Frameworks, wie Xamarin \cite{xamarin_inc._xamarin_2016} oder React Native, vereinen die Vorteile dieser drei Kategorien. Wie bei den hybriden und Web-Apps werden die Apps in einer plattformunabhängig entwickelt. Dazu wird entweder in C\# (Xamarin) oder JavaScript (React Native) benutzt. Die dabei verwendeten Komponenten werden jedoch auf native Komponenten abgebildet, sodass das Endprodukt komplett nativ läuft.
    \section{React Native}

\subsection{Components}

\subsection{Props}

\subsection{State}

\subsection{Styling}
% hier mehr auf farben etc, nicht so sehr auf Layouting -> siehe unten

\subsection{Layouting}
% Flexbox

\subsection{Plattform-Spezialisierung}
https://facebook.github.io/react-native/docs/platform-specific-code.html

\subsection{Integration nativer Module}
% Linking? 


    \section{Projektaufbau}

\subsection{Entwicklungsumgebung}
React Native lässt sich in jedem Texteditor entwickeln. Es gibt aber auch spezielle Entwicklungsumgebungen wie \textit{Nuclide}, welches ein Package für den Texteditor \textit{Atom} ist. Für die Entwicklung mit JavaScript bringt dies jedoch kaum Vorteile und macht Atom zudem spürbar langsamer. Deco ist eine weitere mögliche Entwicklungsumgebung. Diese beinhaltet allerdings nur wenige Features und keine ESLint Integration, dies macht eine qualitative Entwicklung einer App schwierig. Die besten Erfahrungen konnten mit Atom in Kombination mit ESLint und weiteren React Native Packages gemacht werden. 

\subsection{Verzeichnisstruktur}
Mit dem Befehlt \textit{react-native init} lässt sich leicht ein default React Native Projekt generieren, dies erleichtert den Einstieg in ein Projekt. Wie Listing \ref{lst:idirectory_structure} zeigt, besitzt dieses schon eine Reihe von vorgegebenen Dateien und Ordnern. Die Ordner \textit{android} und \textit{ios} beinhalten dabei all den nativen Code der jeweiligen Plattform. Im \textit{ios}-Ordner befindet sich beispielsweise das Xcode-Projekt der App \cite{carli_project_2016}. Passend zu den Ordnern gibt es jeweils eine Index-Datei für Android und eine für iOS. \textit{index.ios.js} ist der Ausgangspunkt für die iOS-App beziehungsweise \textit{index.android.js} für die Android-App. In diesen Dateien wird die App registriert, wie in Kapitel \ref{component} dargestellt. Da die beiden Dateien zu Beginn überwiegend gleichen Code beinhalten, lohnt es sich diesen in eine Datei auszulagern, die schließlich importiert wird. Die eigentliche App wird im Ordner \textit{App} entwickelt. Entwickelte Tests werden im Ordner \textit{\_\_tests\_\_} gespeichert. Hinzukommen noch ein \textit{node\_modules}-Ordner, in der die mittels des \textit{Node Package Manager} installierten Bibliotheken gespeichert werden. Weiterhin liegen verschiedene Konfigurationsdateien im Projektordner. Diese werden weitestgehend von React Native selbst verwaltet. Hervorzuheben sind die \textit{package.json}, welche die installierten Bibliotheken und deren Versionen auflistet, sowie die Konfiguration des JavaScript-Compilers \textit{Babel .babelrc}, welche zum Beispiel beim Hinzufügen von Sprachfeatures, wie speziellen Annotationen, angepasst werden muss. 

  \lstdefinestyle{tree}{
    literate=
    {├}{{\smash{\raisebox{-1ex}{\rule{1pt}{\baselineskip}}}\raisebox{0.5ex}{\rule{1ex}{1pt}}}}1 
    {─}{{\raisebox{0.5ex}{\rule{1.5ex}{1pt}}}}1 
    {└}{{\smash{\raisebox{0.5ex}{\rule{1pt}{\dimexpr\baselineskip-1.5ex}}}\raisebox{0.5ex}{\rule{1ex}{1pt}}}}1 
}
        
\begin{lstlisting}[style=tree]
.
├── .babelrc
├── .buckconfig
├── .flowconfig
├── .watchmanconfig
└── __tests__/
├── android/
├── app/
├── index.android.js
├── index.ios.js
├── ios/
├── node_modules/
└── package.json
    \end{lstlisting}
\vspace{-0.5 cm}
\begin{listing}[H]
    \caption{Initiale Verzeichnisstruktur eines React Native Projekts}
    \label{lst:idirectory_structure}
\end{listing}




\subsection{Ausführen der Anwendung}
Vorausgesetzt Node und React Native wurden entsprechend der online-verfügbaren Anleitung installiert, müssen weiterhin eventuell benötigte Abhängigkeiten, welche in der \textit{package.json} hinterlegt wurden, installiert und gelinkt werden. Soll die Anwendung auf dem Android-Emulator ausgeführt werden, so muss dieser schon vorher gestartet werden.
\begin{listing}[H]
    \begin{minted}{bash}
npm install
react-native link
    \end{minted}
    \caption{Installieren der in der package.json hinterlegten Abhängigkeiten}
    \label{lst:install_dependencies}
\end{listing}

Um die App anschließend auszuführen wird entweder \textit{react-native run-ios} oder \textit{react-native run-android} ausgeführt. Nun wird der Code kompiliert und die App auf dem Simulator gestartet. Werden Änderungen am Projekt vorgenommen, kann die App im Simulator mittels \textit{CMD + R} (iOS) beziehungsweise \textit{RR} (Android) neugeladen werden. Aktiviert man im Developermenü die Option \textit{Hot Reloading} werden neue Quellcodeversionen automatisch geladen ohne aktuellen Zustand zu verlieren. Weiterhin kann die App auch mittels der nativen Umgebungen gestartet werden. Für iOS kann beispielsweise die Datei \textit{<Projekt>/ios/<Projekt>.xcodeproj} mit Xcode geöffnet und ausgeführt werden.

\subsection{Debugging}
Debuggen geht zum einen über die Logs der Simulatoren. Diese können entweder über den Simulator betrachtet oder mittels \textit{react-native log-ios} oder \textit{react-native log-android} betrachtet werden. Zum anderen kann die App mittels der Chrome Developer Tools debugged werden. Dazu öffnet man mittels \textit{CMD + D} (iOS) beziehungsweise \textit{CMD + M} (Android) das Developermenü auf dem Simulator und aktiviert die Option \textit{Debug JS Remotely}. Indem sich nun öffnenden Chrome-Tab können nun die Ausgaben der App in der Konsole verfolgt werden. Es ist auch möglich im Code Breakpoints mittels der Direktive \textit{debugger;} zu setzen, um die Ausführung der App an gewünschten Punkten zu pausieren.  Außerdem kann in der Developerkonsole die Option \textit{Pause On Caught Exceptions} für einfacheres Debugging aktiviert werden.
    \section{Bibliotheken}

\subsection{Einbinden von Bibliotheken}
React Native konzentriert sich auf einige wenige Kernfunktionen. Zur Erweiterung der Funktionalität müssen dem Projekt daher entweder geeignete Bibliotheken hinzugefügt werden oder Code selbst geschrieben werden. Insbesondere falls es für eine gewünschte plattformspezifische Funktionalität noch kein Gegenstück in React Native gibt, muss mittels nativem Code eine Bridge zwischen beidem geschrieben werden. React Native besitzt allerdings eine sehr engagierte Community, sodass dies nur selten notwendig wird.\\
Die Bibliotheken können mittels des Node Package Managers installiert werden. Es bietet sich an die Installation mit dem Flag \textit{-{}-save} durchzuführen, sodass die Bibliothek und deren Version im package.json-File hinterlegt werden. Beinhalten Bibliotheken Funktionalitäten, die mittels nativem Code, also Objective-C, Swift oder Java, umgesetzt sind, müssen die Bibliotheken außerdem noch gelinkt werden. Diese zwei Schritte reichen meist um Bibliotheken dem Projekt hinzuzufügen. Dennoch sollten die Beschreibungen der Libraries auf entsprechende weitere Schritte hin überprüft werden, wie beispielsweise das Hinzufügen von Berechtigungen. 

\begin{listing}[H]
    \begin{minted}{bash}
npm install --save <bibliothek>

# Falls  Bibliothek mit nativem Code
react-native link
    \end{minted}
    \caption{Installation und Linken einer Bibliothek}
    \label{lst:lib_install}
\end{listing}

\subsection{Navigation}
Mobile Apps beinhalten im Regelfall mehrere Ansichten, zwischen denen der Nutzer durch Aktionen navigieren kann. Befindet sich der Nutzer beispielsweise in einer Listenansicht wird er über einen Klick über einen Button \textit{Hinzufügen} auf eine Sicht weiterleitet, in der er einen neuen Listeneintrag hinzufügen kann. React Native stellt zu diesem Zwecks den \textit{Navigator} und \textit{NavigatorIOS} bereit. Letzter kann jedoch ausschließlich für die iOS-Plattform genutzt werden. Der Navigator wird als oberste Komponente definiert, der alle anderen Komponenten beinhaltet. 
\begin{listing}[H]
    \begin{minted}{js}
class App extends React.Component {
    render() {
        const routes = [
            {title: 'First Scene', index: 0},
            {title: 'Second Scene', index: 1},
        ];
        return (
            <Navigator
                initialRoute={routes[0]}
                initialRouteStack={routes}
                renderScene={_renderScene}/>
    }
    
    _renderScene(route, navigitator) {
         if (route.index === 0) {
             return <Scene1 title={route.title} />;
        } else {
            return <Scene2 title={route.title} />;
        }
    }
}
    \end{minted}
    \caption{Beispiel der Navigation mittels des Navigators}
    \label{lst:navigator}
\end{listing}

Wie in Listing \ref{lst:navigator} zu sehen, wird die ganze Navigation von der Top-Level-Komponente aus gesteuert. Dadurch werden Komponenten-spezifische Werte, wie Titel, monolitisch in dieser Komponente gebündet. Außerdem wird die Render-Funktion schon bei wenigen Seiten durch if-else-Kaskaden unübersichtlich. Dieses Problem wird durch Faktoren, wie unterschiedliche Navigationsleisten, noch vergrößert. Daher empfiehlt sich die Verwendung der Bibliothek \textit{ExNavigation} \cite{exNavigation_2016}). Aufgrund der Nützlichkeit und Popularität befindet sich die Bibliothek zur Zeit in der Aufnahme in die \textit{reactjs}-Organisation, wofür sie in \textit{react-navigation} unbenannt wird und zahlreiche Änderungen statt finden. ExNavigation bietet zahlreiche Features, wie Tabs oder Drawer-Navigation, an und erlaubt dabei die Änderungen des Navigationszustandes aus den Kind-Komponenten heraus. Auch wenn das nachfolgende Beispiel nur einen kurzen Einblick in die umfangreichen Funktionen der Bibliothek gewährt, so zeigt doch die Einfachheit und Modularität.

\begin{listing}[H]
    \begin{minted}{js}
 const Router = createRouter(() => ({
     scene1: () => Scene1,
     scene2:() => Scene2
 }));
 
 class App extends React.Component {
     render() {
         return (
             <NavigationProvider router={Router}>
                 <StackNavigation initialRoute={Router.getRoute('home')} />
             </NavigationProvider>
         );
     }
 }
 
 class Scene2 extends React.Component {
     static route = {
         navigationBar: {
             title: 'First Scene',
         }
     }
     // Restliche View
 }
 
 class Scene2 extends React.Component {
     static route = {
         navigationBar: {
             title: 'Second Scene',
         }
     }
    // Restliche View
 }
    \end{minted}
    \caption{Listing \ref{lst:navigator} mit ExNavigation umgesetzt}
    \label{lst:ex_navigation}
\end{listing}

\subsection{Datenbank}
Die React Native Dokumentation führt zur dauerhaften Persistenzierung von Daten die Klasse \textit{AsyncStorage} an, welche es erlaubt Daten im Key-Value-Format zu speichern. Facebook weist dabei darauf hin, \textit{AsyncStorage} nicht direkt zu verwenden, sondern allenfalls indirekt durch eine Bibliothek \cite{AsyncStorage_2017}. Eine weitere Alternative ist die Bibliothek \textit{react-native-sqlite-storage} \cite{react-native-sqlite-storage_2016}, wobei der Nutzer sich selbst um die entsprechenden SQL-Befehle kümmern muss. Empfehlenswert ist daher die Datenbank \textit{Realm}, welche für Java, Objective-C, Swift, Xamarin und eben React Native verfügbar ist und von Firmen wie Google, Amazon oder Intel genutzt wird \cite{Realm_2016}. Realm bietet eine einfache Anfragesprache, Verschlüsselung sowie Benachrichtigungen bei Datenänderungen. Große Stärke ist weiterhin die Performance, welche laut frei verfügbarem Benchmark SQLite sowie AsyncStorage bei weitem übertrifft \cite{Introducing_Realm_React_Native_2016}. 

\begin{listing}[H]
    \begin{minted}{js}
const Dog = {
    name: 'Dog',
    properties: {
        name: 'string',
        age: 'int',
    }
};

let realm = new Realm({schema: [Dog]});

// Speichern eines Hundes
realm.write(() => {
    realm.create('Dog', { name: 'Rex', age: 3 });
});

// Auslesen aller Hunde, die älter als 8 sind
let r = realm.objects('Dog').filtered('age < 8');
    \end{minted}
    \caption{Beispielhafte Realm-Anfragen}
    \label{lst:realm}
\end{listing}

\subsection{Internationalisierung}
Die Anzeige lässt sich dank der Bibliothek \textit{react-native-i18n} einfach internationalisieren \cite{react-native-i18n_2017}. Dazu werden englischen Formulieren sowie deren Übersetzung anhand eines eindeutigen Keys identifiziert. Sofern die Option \textins{Fallback} aktiviert ist, wird bei Fehlen einer entsprechenden Übersetzung automatisch auf die englische Version zurückgegriffen. Standardmäßig werden Texte automatisch anhand der festgelegten Gerätesprache ausgewählt.

\begin{listing}[H]
    \begin{minted}{js}
class Example extends React.Component {
    render () {
        return (
            <Text>{I18n.t('hello_word')}</Text>
        );
    }
}

I18n.fallbacks = true

I18n.translations = {
    en: {
        hello_word: 'Hello world!'
    },
    de: {
        hello_word: 'Hallo Welt!'
    }
}
    \end{minted}
    \caption{Internationalisierung mittels \textit{i18n}}
    \label{lst:internationalizing}
\end{listing}


%TODO
% \subsection{Testing}
   % \section{Waterings of Things Start}

	\renewcommand{\listingscaption}{Code Example}

\begin{listing}[H]
	\begin{minted}{python}
class UserViewSet(viewsets.ModelViewSet):
"""
API endpoint that allows users to be viewed or edited.
"""
queryset = User.objects.all().order_by('-date_joined')
serializer_class = UserSerializer
	\end{minted}
\caption{Views}
\end{listing}

\begin{listing}[H]
    \begin{minted}{js}
var variable = "Beispieltext";
alert(typeof variable);
    \end{minted}
    \caption{JavaScript}
\end{listing}

\begin{listing}[H]
    \begin{minted}{cpp}
void printNames(const std::vector<std::string>& names)
{
for (const std::string& singleName : names)
{
std::cout << singleName << std::endl;
}
}
    \end{minted}
    \caption{C++}
\end{listing}
   	\section{Résumé}
Mit dem Einsatz von WateringOfThings kann der Nutzer nun seine Pflanzen auch längere Zeit über alleine lassen, ohne auf eine ideale Bewässerung der Pflanzen verzichten zu müssen. Das Projekt stellt dabei eine Verknüpfung von Webtechnologien und Internet of Things dar. Dabei wurden zum einen eine native Smartphone-Applikation geschrieben, die sowohl auf Android als auch iOS lauffähig ist. Zum anderen wurde mithilfe von Django ein Web-API umgesetzt, die die Kommunikation mit den Hardwarekomponenten übernimmt sowie die gesammelten Daten zentral verwaltet und bereitstellt. Darüber hinaus wurde eine eigene Hardwarekomponente, beginnend von der Auswahl der benötigten Teile bis hin zur Umsetzung eines eigenen Platinenlayouts, konzeptioniert und in die Gesamtarchitektur mittels des leichtgewichtigen MQTT-Protokolls eingebunden. Im Zusammenspiel setzen diese drei Komponenten alle in Kapitel \ref{sec:usecases} beschriebenen Use Cases um. Aufgrund des Nutzen von WateringOfThings wird eine Fortführung des Projektes angestrebt. Dafür kommen folgende Features in Frage:
\begin{itemize}
    \item Hinzufügen eines weiteren Sensors, um den Nutzer über Fehlfunktionen zu benachrichtigen
    \item Weitere Sensoren (Licht, Temperatur), um die Bewässerung zu automatisieren
    \item Bündelung der Hardware in einem sich selbst bewässernden Blumentopf
    \item Push-Notifications, um den Nutzer über zu trockene Pflanzen zu informieren
    \item Möglichkeit mehrere Hardwarekomponenten pro App zu verwalten
    \item Nutzertests um die UI nutzerfreundlicher zu gestalten
    \item Visualisierung des Verlaufs der Feuchtigkeit einer Pflanze
    \item Hinzufügen weiterer Sensoren, um den Düngebedarf der Blumenerde festzustellen
\end{itemize}


    \newpage
    \appendix

        \section{Anhang}\label{anhang}

    \clearpage
	\printbibliography %Literaturverzeichnis
\end{document}
